\documentclass[10pt]{article}
\usepackage[a4paper,margin=2cm]{geometry}
\usepackage[T1]{fontenc}
\usepackage{lmodern}
\usepackage{microtype}
\usepackage{booktabs}
\usepackage{tabularx}
\usepackage{xltabular}
\usepackage{array}
\usepackage{xcolor}
\usepackage{enumitem}
\usepackage[hidelinks]{hyperref}

\definecolor{pending}{HTML}{B3261E}
\newcommand{\pending}[1]{\textcolor{pending}{\textbf{Pending:} #1}}
\newcolumntype{L}{>{\raggedright\arraybackslash}X}
\newcolumntype{R}[1]{>{\raggedright\arraybackslash}p{#1}}
\renewcommand{\arraystretch}{1.25}
\setlist{nosep,leftmargin=*}

\title{Response to the Reviewers\\[2pt]\large CICT-2026 --- ``Web2Music: Page-Conditioned Music Generation with a Disclosure-Aware Classifier Cascade''}
\author{}
\date{}

\begin{document}
\maketitle
\vspace{-3em}

\noindent We thank both reviewers for a careful reading and for recommending acceptance after
minor revision. Each comment is answered below with what changed in the manuscript and
where. Locations refer to section and table titles in the revised paper. Every number
quoted here is produced by a script in the released repository from a result file that is
committed alongside it.

\medskip
\noindent Four requests could not be met in the revision window: we ran \textbf{no human listening
study}, we could \textbf{not recruit independent annotators}, we did \textbf{not build a
held-out test set}, and we did \textbf{not evaluate the sensitive-page detector on real
crisis pages}. For each we say what we did instead and how the claims were narrowed; the
paper states all four as limitations and contains no placeholder for any of them.

\section*{Summary of the main changes}
\begin{itemize}
\item \textbf{Real end-to-end latency on a consumer GPU.} The built extension was measured
  with the real embedding model, classifier and generator (110 page loads per setup).
  Generation was measured on an RTX~5060 laptop GPU (33 cold requests per clip length), on the
  same machine with the GPU disabled, and on an Apple~M4.
\item \textbf{How often generated audio is heard.} A browsing simulation now uses measured
  generation times and the cache hit rate measured on real page profiles. With a cold cache,
  page-level audio is heard on 52\% of visits on the GPU, 48\% on the M4 and 34\% on the CPU
  (15\,s clips, median dwell 30\,s).
\item \textbf{Objective conditioning fidelity}, in place of a listening study: 702 clips for
  234 pages compared against shuffled-mood controls on three scorers. The effect is
  significant but small, and the paper says so.
\item \textbf{Loops.} A fresh seeded study of 220 clips at the shipped crossfade, objective
  gapless verification (30 of 30), and seam size relative to each clip's own variation.
\item \textbf{Generator comparison} against MusicGen-medium and AudioLDM2-music.
\item \textbf{Network-level privacy audit} of the built extension, and a switch that
  makes the fully local configuration (A6) send no LLM request at all.
\item \textbf{Prompt ablation} rerun with fixed seeds and paired tests. The earlier ``template wins''
  reading did not reproduce and was replaced by the null result.
\end{itemize}

\section*{Requests we could not meet in full}

\paragraph{Listening study (R3 W1/R1; R4 \#2).}
We could not run a human listening study within the revision period, so we:
(i) removed every claim that the music is perceived as fitting the page;
(ii) added objective tests of conditioning fidelity; and
(iii) separated \emph{gapless} playback, which we now verify from rendered output, from
\emph{perceptual} seamlessness, which we do not claim.

The fidelity tests use audio valence/arousal regressors, a mood tagger and CLAP retrieval.
Each compares page-matched against shuffled-mood clips over 234 pages $\times$ 3 seeds. The
protocol and instruments for a within-subjects listening study are released for follow-up
work.

\paragraph{Independent annotators and a held-out test set (R3 W3/R3; R4 \#1).}
We could not recruit independent annotators, and we did not build the proposed held-out
test set. What we did instead:
\begin{itemize}
\item the paper now says explicitly that the cascade's gates are tuned and scored on the same
  255 development pages, so its macro-F1 values are in-sample and may be optimistic by an
  unmeasured amount; the abstract and conclusion carry the same caveat, and the held-out
  claims were removed;
\item LLMs are kept out of page selection, labelling and judging, and the LLM stage's raw
  answers are cached so every replay is identical;
\item every claim that depends on mood labels is removed, since mood is no longer an
  evaluated target;
\item author annotation, the absence of a held-out set and the small development set
  ($\pm.05$ macro-F1 intervals) are stated as limitations.
\end{itemize}

\section*{Reviewer 3}

\begin{xltabular}{\textwidth}{@{}R{1.25cm}R{3.6cm}LR{3.2cm}@{}}
\toprule
\textbf{\#} & \textbf{Comment} & \textbf{What we did} & \textbf{Where} \\
\midrule
\endhead
W1 / R1 & Fit to the page is never evaluated with listeners &
Narrowed the claim to music \emph{conditioned on} the page, and removed ``suits the page''
wording, including from the abstract. Added objective conditioning-fidelity measures
(702 clips, matched against shuffled-mood controls):
\begin{itemize}
\item matched clips are closer to the intended valence--arousal by 0.050 [0.026, 0.076];
\item the intended mood is in the top-5 tags 7.5 points more often [1.3, 13.8];
\item CLAP page-text similarity is higher by 0.023 [0.014, 0.033];
\item valence and arousal track the targets ($\rho = 0.40$ and $0.51$).
\end{itemize}
The effect is small: CLAP picks a page's own clip among ten in 12.5\% of cases against 10\%
chance. We report it as a tendency. Listening protocol released. &
Abstract; Introduction; Results: Conditioning fidelity (Table ``Conditioning fidelity''); Limitations \\
\midrule
W2 / R2 & ``36\,ms'' is the fallback clip, and users mostly hear the fallback &
Restated the claim as three levels of conditioning: L1 mood-level, L2 profile-level from the
cache, L3 page-level. Measured cold generation on a consumer GPU: a median of 13.3\,s for a
15\,s clip and 24.3\,s for a 28\,s clip (real-time factor 0.87--0.91 end to end, $n=33$ per
length).
Simulated how often each level is heard, with a cold per-session cache and the pre-warm-grid
hit rate measured on real profiles (0 of 235). At 15\,s clips and a median dwell of 30\,s:
\begin{itemize}
\item L3 is heard on 52\% of visits on the GPU, 48\% on the M4 and 34\% on the CPU;
\item L3 is 75\% of listening time on the GPU;
\item the results are stratified by page type, with a dwell sweep from 10 to 120\,s and two
  stated bounds (an assumed 50\% grid hit, and a cache persisting across sessions).
\end{itemize} &
Abstract; Introduction; Results: Generation latency, How often generated audio is heard (Table ``Exposure'', Eq.~4) \\
\midrule
W3 / R3 & Labels come from the authors; $\alpha$ below threshold; mood $\alpha = 0.195$ &
Not resolved (see above). Mood is no longer an evaluated target, the existing agreement
coefficients are reported as they were, and author labelling is stated as a limitation. We
did not collect a new labelled set, so there is no new $\alpha$. &
Results: Evaluation design; Limitations \\
\midrule
W4 & Headline latencies come from stubbed benchmarks &
Re-measured end to end with the real embedding model, classifier and generator: 110 page loads
per setup, all of which became audible, with none timed out.
\begin{itemize}
\item With the LLM stage off, the first sound is heard at a median of 7.7\,s (p95 12.1\,s),
  and the generated clip takes over by p95 37.6\,s.
\item Descriptor cost with the real, warm embedding model: median 360\,ms (p95 519\,ms) on 166
  live sites.
\end{itemize}
The stubbed figures remain only where labelled as such. &
Abstract; Introduction; Results: Descriptor cost and classification, Generation latency (Table ``Latency'') \\
\midrule
R4 (clarity) & Garbled sentences; ``bart-large-mini''; dense notation &
Rewrote the abstract, introduction and the architecture sections in plain language. The
models are named \texttt{bart-large-mnli} and \texttt{nli-deberta-v3-xsmall} throughout. Four
numbered equations (gate rule, disclosure, seam cost, exposure) are each explained in words,
and terms (macro-F1, NLI, off-device rate, RTF, LUFS, $\alpha$) are defined at first use. &
Throughout \\
\bottomrule
\end{xltabular}

\section*{Reviewer 4}

\begin{xltabular}{\textwidth}{@{}R{1.25cm}R{3.6cm}LR{3.2cm}@{}}
\toprule
\textbf{\#} & \textbf{Comment} & \textbf{What we did} & \textbf{Where} \\
\midrule
\endhead
1 & Independent held-out corpus and annotators; tune thresholds on a separate set &
Not done. We did not build a held-out set or recruit independent annotators. The paper
now states that the gates are tuned on the pages they are scored on, labels the results
in-sample, and lists both gaps as limitations. &
Abstract; Results: Evaluation design, Classification and disclosure; Conclusion; Limitations \\
\midrule
2 & Listening study for fit and gapless loops &
A listening study was not possible. Instead:
\begin{itemize}
\item \textbf{Gapless playback, verified objectively:} rendering clips through the browser's
  decoder for 10 loop cycles, 30 of 30 show no inserted silence and are sample-identical from
  cycle to cycle.
\item \textbf{Seam size relative to each clip's own variation:} the spectral-centroid change at
  the seam is within the clip's 95th percentile in 91.4\% [87.7, 95.0] of 220 clips. The
  energy share (100\%) is reported as circular, because the selector minimises it.
\end{itemize}
``Seamless'' and ``inaudible'' are not claimed. &
Results: Loop quality \\
\midrule
3 & Improve and re-measure the sensitive detector, including the on-device model &
Partly done. The shipped check now adds an ``unrelated'' hypothesis (three-way entailment)
and runs on-device. We did not evaluate it on real crisis pages or on benign pages, so its
recall and false-alarm rate are unmeasured, and the paper says so. What we could measure is
the network audit, which shows that no page text leaves the device in the fully local
configuration. &
Results: Restraint (Table ``Sensitive-page detection''); Limitations \\
\midrule
4 & Real end-to-end latency; enough runs for p95; CPU-only case &
See R3 W4.
\begin{itemize}
\item \textbf{Enough runs for p95:} each cold-generation cell has 33 requests, and each
  end-to-end setup 110 loads.
\item \textbf{CPU-only:} the GPU host with the GPU disabled generates a 10\,s clip in a median
  of 38.5\,s and a 15\,s clip in 62.2\,s (12 per length, none timed out). The simulation shows
  page-level audio on 34\% of visits in this configuration.
\end{itemize} &
Results: Generation latency (Table ``Latency'', Fig.\ ``Cold generation''); How often generated audio is heard \\
\midrule
5 & Seeds and intervals for prompt and loop studies; compare another generator &
Fixed seed lists, paired tests and bootstrap CIs:
\begin{itemize}
\item \textbf{Loop study:} 220 clips, 176 scoreable, with the two exclusion reasons stated.
  The share of seams above 1\,dB falls from 72.7\% to 8.5\% (McNemar $p = 2\times10^{-34}$).
\item \textbf{Prompt study:} 165 clips. The earlier difference does not reproduce (engineered
  minus template 0.00\,dB [$-0.10$, 0.04], $p = 0.88$), and the paper now reports a null
  result.
\item \textbf{Generator comparison:} MusicGen-small, MusicGen-medium and AudioLDM2-music, 55
  clips each, reporting RTF, seam, loop retention, CLAP-space FAD and Audiobox quality.
\end{itemize} &
Results: Loop quality; Conditioning fidelity and generator choice (generator comparison and prompt ablation, in text) \\
\midrule
6 & Rewrite abstract and intro; define terms; release code and data &
Rewritten, and terms defined. Code, URLs, result files and the raw LLM responses are released
(page text withheld for copyright); the archive DOI is added when the repository is released. &
Abstract; Introduction; repository release \\
\midrule
7 & Mathematical modelling, annotation reliability, privacy validation, perception, clarity &
\textbf{Modelling and annotation:} four numbered equations (gate rule, disclosure, seam cost,
exposure), and $\alpha$ reported on the development set.

\textbf{Privacy validation:} a network-level audit of the built extension.
\begin{itemize}
\item In the fully local configuration (A6) no request body carries page text.
\item In the default (A6b) page text goes only to the LLM proxy.
\item No telemetry leaves the browser, and every request to the generation service carries
  only the music profile.
\item URL hashes in telemetry now use a keyed HMAC.
\end{itemize}
\textbf{Perception:} handled as above. &
System architecture; Results: Evaluation design, Restraint (network audit) \\
\midrule
W & Limited comparison with alternatives &
Generator comparison added. The cascade table's seven configurations, from keyword-only to
LLM-only, are the baselines and are now named as such. &
Table ``Cascade configurations''; Results: Conditioning fidelity and generator choice \\
\bottomrule
\end{xltabular}

\section*{Other corrections found during the revision}
\begin{itemize}
\item \textbf{The shipped build now matches the evaluated system.} The on-device zero-shot
  middle stage and the measured-seam loop selector are the defaults.
\item \textbf{Crossfade.} The paper previously described a 50\,ms default crossfade; the
  shipped default is 250\,ms, and the text now says so.
\item \textbf{The fully local configuration (A6) previously still attempted LLM requests.} It
  sent them to the local proxy address and they failed when the proxy was down. The network
  audit exposed this. A6 is now an explicit setting under which no LLM request is made, and
  the audit and the A6 latency run were repeated with it.
\item \textbf{The exposure simulation earlier assumed that half of all pages hit the pre-warm
  cache.} Replaying real page profiles gives 0 of 235, because cache keys include the exact
  style string. The simulation now uses the measured value, and the assumed value is reported
  only as a bound.
\item \textbf{Real-time factors.} They are now reported end to end (0.87--0.91), with the
  generator-only figure (0.83--0.85) stated separately.
\item \textbf{The descriptor-cost benchmark previously timed the embedding model's download
  and load on every page.} It now times per-page inference with the model loaded once, as in
  the extension.
\item \textbf{Batching.} Concurrent generation requests were measured to give no throughput
  gain on the GPU (speedup 1.01--1.02$\times$ at 2--8 concurrent requests, five repeats), and
  the paper says so.
\end{itemize}

\end{document}
